\documentclass[
  12pt,
  a4paper,
  oneside,
  final
]{report}

% ─── Encoding & Language ────────────────────────────────────────────────────
\usepackage[utf8]{inputenc}
\usepackage[T1]{fontenc}
\usepackage[indonesian]{babel}
\usepackage{indentfirst}
\usepackage{setspace}
\onehalfspacing

% ─── Page Layout ────────────────────────────────────────────────────────────
\usepackage{geometry}
\geometry{
  left   = 3.0cm,
  right  = 2.5cm,
  top    = 3.0cm,
  bottom = 2.5cm,
  headheight = 1.2cm,
  headsep  = 0.8cm,
  footskip = 1.2cm
}

% ─── Fonts ───────────────────────────────────────────────────────────────────
\usepackage{times}       % Times New Roman

% ─── Colors ─────────────────────────────────────────────────────────────────
\usepackage[table]{xcolor}
\definecolor{navy}{RGB}{30, 58, 95}
\definecolor{lightblue}{RGB}{219, 229, 241}
\definecolor{accent}{RGB}{0, 100, 160}
\definecolor{lightgray}{RGB}{245, 245, 245}
\definecolor{darkgray}{RGB}{80, 80, 80}
\definecolor{superadmin}{RGB}{139, 0, 0}
\definecolor{admin}{RGB}{0, 100, 0}
\definecolor{peminjam}{RGB}{0, 60, 130}

% ─── Hyperlinks & Bookmarks ─────────────────────────────────────────────────
\usepackage[
  colorlinks=true,
  linkcolor=navy,
  urlcolor=navy,
  citecolor=navy,
  anchorcolor=navy,
  pdfdisplaydoctitle=true,
  pdftitle={Dokumentasi SIPP-BMN},
  pdfauthor={Tim Pengembang},
  pdfsubject={Sistem Informasi Peminjaman \& Pengembalian Barang Milik Negara}
]{hyperref}
\usepackage{bookmark}

% ─── Tables ─────────────────────────────────────────────────────────────────
\usepackage{longtable}
\usepackage{booktabs}
\usepackage{colortbl}
\usepackage{tabularx}
\usepackage{multirow}
\usepackage{array}
\newcolumntype{L}{>{\raggedright\arraybackslash}X}
\newcolumntype{C}{>{\centering\arraybackslash}X}
\newcolumntype{R}{>{\raggedleft\arraybackslash}X}

% ─── Headers & Footers ──────────────────────────────────────────────────────
\usepackage{fancyhdr}
\pagestyle{fancy}
\fancyhf{}
\fancyhead[L]{\fontsize{10}{12}\selectfont\textbf{Dokumentasi SIPP-BMN}}
\fancyhead[R]{\fontsize{10}{12}\selectfont\textbf{\thepage}}
\fancyfoot[C]{\fontsize{9}{11}\selectfont\thepage}
\fancypagestyle{plain}{
  \fancyhf{}
  \fancyfoot[C]{\fontsize{9}{11}\selectfont\thepage}
}
\setlength{\headrulewidth}{0.5pt}
\setlength{\footrulewidth}{0pt}

% ─── Section Formatting ──────────────────────────────────────────────────────
\usepackage{titlesec}
\titleformat{\chapter}[display]
  {\normalfont\fontsize{16}{18}\selectfont\bfseries\color{navy}}
  {\chaptername\ \thechapter}{0.5em}{}
\titleformat{\section}
  {\normalfont\fontsize{14}{16}\selectfont\bfseries\color{navy}}
  {\thesection}{0.5em}{}
\titleformat{\subsection}
  {\normalfont\fontsize{12}{14}\selectfont\bfseries\color{navy}}
  {\thesubsection}{0.5em}{}
\titleformat{\paragraph}
  {\normalfont\fontsize{11}{13}\selectfont\bfseries\color{darkgray}}
  {}{0em}{}
\titlespacing*{\chapter}{0pt}{0pt}{20pt}
\titlespacing*{\section}{0pt}{18pt}{10pt}
\titlespacing*{\subsection}{0pt}{12pt}{8pt}
\titlespacing*{\paragraph}{0pt}{8pt}{6pt}

% ─── Captions ───────────────────────────────────────────────────────────────
\usepackage[
  tableposition=top,
  figureposition=bottom
]{caption}
\captionsetup{
  font=small,
  labelfont=bf,
  labelsep=colon,
  justification=raggedright
}
\captionsetup[table]{
  font={small,bf},
  skip=6pt,
  aboveskip=8pt
}

% ─── Lists ─────────────────────────────────────────────────────────────────
\usepackage{enumitem}
\setlist{
  topsep=4pt,
  itemsep=3pt,
  parsep=2pt,
  labelsep=0.4cm
}
\setlist[itemize,1]{label={\textbullet}}
\setlist[itemize,2]{label={--}}
\setlist[itemize,3]{label={\textendash}}

% ─── Graphics ───────────────────────────────────────────────────────────────
\usepackage{graphicx}
\graphicspath{{./assets/}}

% ─── Misc ───────────────────────────────────────────────────────────────────
\usepackage{pdfpages}
\usepackage{amsmath}
\usepackage{amssymb}
\usepackage{etoolbox}
\usepackage{ifthen}
\usepackage{xstring}

% ─── Document Begins ────────────────────────────────────────────────────────
\begin{document}

% ─── Cover Page ────────────────────────────────────────────────────────────
\pagenumbering{alph}
\thispagestyle{empty}

\begin{center}
  \vspace*{1.5cm}

  % Logo placeholder
  \IfFileExists{./assets/logo.png}{%
    \includegraphics[width=5cm]{./assets/logo.png}\\[1.5cm]
  }{%
    % Decorative top line
    \vspace*{1.5cm}
  }

  \LARGE\textbf{\textls{DOKUMENTASI SISTEM}}\\[0.3cm]
  \LARGE\textbf{SIPP-BMN}\\[0.5cm]

  \HRule\\[0.3cm]

  {\LARGE\textbf{Sistem Informasi Peminjaman \&}}
  \\[0.3cm]
  {\LARGE\textbf{Pengembalian Barang Milik Negara}}\\[0.5cm]

  \HRule\\[1.5cm]

  \begin{tabular}{rl}
    \textbf{Versi}      & : 1.0.0 \\[4pt]
    \textbf{Tanggal}     & : \today \\[4pt]
    \textbf{Status}      & : Final
  \end{tabular}\\[2cm]

  \vfill

  {\large\textbf{Tim Pengembang SIPP-BMN}}

  \vspace{1cm}
\end{center}

\newpage
\thispagestyle{empty}
~
\newpage

\pagenumbering{roman}
\pagestyle{fancy}

% ─── Table of Contents ──────────────────────────────────────────────────────
\tableofcontents
\newpage
\listoffigures
\addcontentsline{toc}{chapter}{Daftar Gambar}
\newpage
\listoftables
\addcontentsline{toc}{chapter}{Daftar Tabel}
\newpage

\pagenumbering{arabic}
\pagestyle{fancy}

% ─── Chapter 1: Pendahuluan ─────────────────────────────────────────────────
\chapter{Pendahuluan}

\section{Latar Belakang}
\index{latar belakang}
Perkembangan teknologi informasi di lingkungan instansi pemerintah menuntut
adanya sistem yang efisien, transparan, dan akuntabel dalam pengelolaan
Barang Milik Negara (BMN). Selama ini, proses peminjaman dan pengembalian
barang milik negara masih mengandalkan dokumen fisik dan pencatatan manual
yang memiliki berbagai keterbatasan, antara lain:
\begin{itemize}
  \item sulitnya menjaga konsistensi data antar unit kerja;
  \item tidak adanya informasi stok secara {\it real-time};
  \item risiko kehilangan atau kerusakan dokumen;
  \item sulitnya melakukan pelacakan dan audit terhadap transaksi
        peminjaman; dan
  \item proses persetujuan yang memakan waktu lama karena ketergantungan
        pada kehadiran fisik pihak yang berwenang.
\end{itemize}

Untuk menjawab tantangan tersebut, dikembangkanlah {\bfseries SIPP-BMN}
({\it Sistem Informasi Peminjaman \& Pengembalian Barang Milik Negara}),
sebuah aplikasi berbasis web yang menyediakan pengelolaan peminjaman dan
pengembalian BMN secara {\it end-to-end}, mulai dari pengajuan oleh
peminjam, persetujuan oleh administrator, hingga generasi dokumen resmi
berupa Surat Pernyataan Peminjaman dan Surat Pengembalian secara otomatis.

\section{Tujuan Dokumentasi}
\index{tujuan dokumentasi}
Dokumentasi ini dirancang untuk:
\begin{itemize}
  \item memberikan gambaran menyeluruh tentang arsitektur dan fitur sistem;
  \item menjadi panduan bagi setiap jenis pengguna ({\it role}) dalam
        mengoperasikan aplikasi sesuai dengan hak akses masing-masing;
  \item mendokumentasikan entitas basis data dan alur kerja utama;
  \item memudahkan proses pengembangan, pemeliharaan, dan transisi
        pengetahuan ({\it knowledge transfer}) bagi tim teknis; dan
  \item Serve sebagai acuan dalam proses pelatihan ({\it training}) dan
        {\it user acceptance testing} (UAT).
\end{itemize}

\section{Ruang Lingkup}
\index{ruang lingkup}
Dokumentasi ini mencakup tiga kategori utama, yaitu:
\begin{enumerate}
  \item {\bfseries Fitur dan Halaman} --- daftar lengkap halaman yang
        tersedia untuk setiap peran pengguna.
  \item {\bfseries Alur Kerja} --- penjelasan proses bisnis dari awal
        pengajuan hingga pengembalian barang.
  \item {\bfseries Arsitektur Sistem} --- entitas basis data, teknologi
        yang digunakan, serta mekanisme keamanan.
\end{enumerate}

Dokumentasi ini {\bfseries tidak} mencakup panduan instalasi lokal
({\it local development setup}), konfigurasi variabel lingkungan
({\it environment variables}) pada level infrastruktur, serta kode
sumber program ({\it source code}).

\section{Acuan Dokumen}
\index{acuan dokumen}
Dokumen ini disusun berdasarkan analisis kode sumber dan struktur basis
data aplikasi SIPP-BMN versi terkini.

% ─── Chapter 2: Gambaran Sistem ─────────────────────────────────────────────
\chapter{Gambaran Sistem}

\section{Deskripsi Sistem}
\index{deskripsi sistem}
SIPP-BMN adalah aplikasi {\it full-stack} berbasis web yang dirancang
untuk mengelola seluruh aspek siklus hidup peminjaman dan pengembalian
Barang Milik Negara. Sistem ini menyediakan fungsi-fungsi utama berikut:

\begin{itemize}
  \item {\bfseries Manajemen Katalog Barang}: Penambahan, pengeditan,
        pencarian, dan penyaringan barang BMN berdasarkan kategori,
        kondisi, lokasi penyimpanan, dan Satuan Kerja (Satker).
  \item {\bfseries Manajemen Peminjaman}: Pengajuan, persetujuan,
        penolakan, penandoveran, dan pengembalian barang dengan
        pelacakan status secara {\it real-time}.
  \item {\bfseries Manajemen Pengguna}: Registrasi, autentikasi,
        otorisasi berbasis peran (RBAC), dan pengelolaan sesi.
  \item {\bfseries Generasi Dokumen Otomatis}: Pembuatan Surat
        Pernyataan Peminjaman dan Surat Pengembalian secara otomatis
        dengan nomor urut, perangko digital ({\it digital stamp}),
        dan tanda tangan digital.
  \item {\bfseries Verifikasi QR Code}: Pembangkitan dan pemindaian
        kode QR untuk verifikasi transaksi pengembalian.
  \item {\bfseries Impor dan Ekspor Data}: Impor dan ekspor data barang,
        pegawai, dan peminjaman dalam format Microsoft Excel.
  \item {\bfseries Sistem Notifikasi}: Pemberitahuan dalam aplikasi
        kepada pengguna terkait status pengajuan dan jadwal pengembalian.
  \item {\bfseries Log Audit}: Pencatatan seluruh aktivitas pengguna
        untuk keperluan audit dan akuntabilitas.
\end{itemize}

\section{Stack Teknologi}
\index{stack teknologi}
Sistem ini dibangun menggunakan teknologi modern yang dipilih berdasarkan
pertimbangan performa, keamanan, skalabilitas, dan kemudahan pemeliharaan.

\begin{table}[!htbp]
  \centering
  \caption{Stack Teknologi SIPP-BMN}
  \label{tab:tech-stack}
  \rowcolors{1}{white}{white}
  \begin{tabularx}{\textwidth}{L L}
    \toprule
    \textbf{Lapisan} & \textbf{Teknologi} \\
    \midrule
    {\bfseries Frontend}  & Next.js 14 (App Router), TypeScript, Tailwind CSS,
                           shadcn/ui (Radix UI), Zustand (state management),
                           Axios, React Hook Form, Zod, date-fns (locale ID),
                           html5-qrcode, lucide-react \\
    {\bfseries Backend}   & Node.js, Express.js, Prisma ORM, PostgreSQL,
                           bcryptjs (password hashing), Multer (file upload),
                           Vercel Blob (penyimpanan file), qrcode,
                           pdf-lib (generasi PDF), xlsx (impor/ekspor Excel),
                           Zod (validasi), helmet (keamanan header),
                           express-rate-limit (rate limiting),
                           nodemailer (email), node-cron (penjadwal) \\
    {\bfseries Penyimpanan} & Vercel Blob (foto, QR code, dokumen PDF) \\
    {\bfseries Autentikasi} & JWT Access Token (60 menit) +
                              Refresh Token via httpOnly cookie (7 hari),
                              CSRF Protection (Double Submit Cookie),
                              Token Versioning, Token Blacklisting \\
    \bottomrule
  \end{tabularx}
\end{table}

\section{Model Arsitektur}
\index{model arsitektur}
SIPP-BMN mengadopsi arsitektur {\it client-server} dengan pola
RESTful API. Arsitektur ini memisahkan concerns antara frontend
dan backend secara jelas, memungkinkan pengembangan dan penskalaan
masing-masing secara independen.

Arsitektur sistem dapat digambarkan sebagai berikut:

\begin{figure}[htbp]
  \centering
  \IfFileExists{./assets/arsitektur.png}%
    {\includegraphics[width=0.85\textwidth]{./assets/arsitektur.png}}%
    {%
      \begin{tabular}{|p{3cm}|p{9cm}|}
        \hline
        \textbf{Lapisan} & \textbf{Komponen} \\
        \hline
        Client & Browser / Next.js 14 (React) \\
        \hline
        API Layer & Express.js (Routing, Middleware, Validasi) \\
        \hline
        Business Logic & Controller, Service, Workflow \\
        \hline
        Data Access & Prisma ORM / PostgreSQL \\
        \hline
        External Services & Vercel Blob, Email Server \\
        \hline
      \end{tabular}
    }
  \caption{Arsitektur Sistem SIPP-BMN}
  \label{fig:arsitektur}
\end{figure}

Klien berkomunikasi dengan backend melalui API REST dengan format
JSON. File statis (foto barang, dokumen PDF, QR code) disimpan di
layanan Vercel Blob dan disajikan melalui URL yang aman.

% ─── Chapter 3: Peran dan Hak Akses ─────────────────────────────────────────
\chapter{Peran dan Hak Akses}

\index{peran pengguna}
\index{hak akses}
\index{role-based access control}
Sistem SIPP-BMN menerapkan model {\it Role-Based Access Control} (RBAC)
dengan tiga peran utama. Setiap akun dapatAssigned lebih dari satu peran
secara bersamaan ({\it multi-role}), yang memungkinkan fleksibilitas
dalam Pendelegasian tugas.

\section{Daftar Peran}
\index{daftar peran}

\begin{longtable}[!htbp]{p{3.5cm} p{4cm} p{6.5cm}}
  \caption{Daftar Peran dalam Sistem}
  \label{tab:roles} \\

  \toprule
  \textbf{Kode Peran} & \textbf{Nama} & \textbf{Deskripsi} \\
  \midrule
  \endfirsthead

  \caption[]{Daftar Peran dalam Sistem (sambungan)} \\
  \toprule
  \textbf{Kode Peran} & \textbf{Nama} & \textbf{Deskripsi} \\
  \midrule
  \endhead

  \bottomrule
  \endfoot

  \rowcolor{lightblue}\textbf{SUPER\_ADMIN} & Super Admin &
  Akses penuh ke seluruh data dan fungsi sistem tanpa batasan
  Satker. Mampu mengelola akun admin, seluruh barang, seluruh
  transaksi, seluruh pengguna, seluruh Satker, dan seluruh log. \\
  \rowcolor{white}\textbf{ADMIN} & Administrator &
  Mengelola barang untuk Satker yang ditugaskan, menyetujui/menolak/
  mengembalikan peminjaman di lingkup Satker-nya, mengelola akun
  peminjam di Satker-nya, serta melihat log aktivitas. \\
  \rowcolor{white}\textbf{PEMINJAM} & Peminjam &
  Mengakses katalog barang, menambahkan barang ke keranjang,
  mengajukan permintaan peminjaman, mengunggah surat pernyataan
  yang ditandatangani, meminta pengembalian, serta melihat
  riwayat pribadi. \\
\end{longtable}

\section{Sistem Multi-Peran}
\index{multi-role system}
\index{role switching}
Sistem ini mendukung {\it multi-role assignment}, di mana satu akun
dapat memiliki kombinasi peran secara bersamaan, misalnya seorang
pengguna sekaligus berperan sebagai {\roleAdmin} dan {\rolePeminjam}.
Pengaturan ini memberikan fleksibilitas operasional yang tinggi
bagiinstansi yang memiliki jumlah pengguna terbatas.

\begin{itemize}
  \item {\bfseries Pemilihan Peran saat Login}: Pengguna dengan beberapa
        peran akan diminta memilih peran yang ingin digunakan saat masuk
        ke sistem.
  \item {\bfseries Pengalihan Peran ({\it Role Switching})}: Pengguna
        dapat beralih antar-peran melalui {\it header switcher} tanpa
        harus keluar dan masuk kembali ke sistem.
  \item {\bfseries Logika Peran Efektif}: Peran {\roleAdmin}
        secara otomatis mencakup kemampuan {\rolePeminjam},
        dan peran {\roleSuperAdmin} mencakup kemampuan {\roleAdmin}
        maupun {\rolePeminjam}.
\end{itemize}

\section{Matriks Hak Akses}
\index{matriks hak akses}

\begin{longtable}[!htbp]{p{5cm} p{1.8cm} p{1.8cm} p{1.8cm}}
  \caption{Matriks Hak Akses per Peran}
  \label{tab:access-matrix} \\

  \toprule
  \multirow{2}{\textbf{Fitur}} &
    \multicolumn{1}{c}{\rotatebox{90}{\textbf{SUPER ADMIN}}} &
    \multicolumn{1}{c}{\rotatebox{90}{\textbf{ADMIN}}} &
    \multicolumn{1}{c}{\rotatebox{90}{\textbf{PEMINJAM}}} \\
  \cmidrule{2-4}
  & \centering \textbf{\roleSuperAdmin} &
    \centering \textbf{\roleAdmin} &
    \centering \textbf{\rolePeminjam} \\
  \midrule
  \endfirsthead

  \caption[]{Matriks Hak Akses per Peran (sambungan)} \\
  \toprule
  \multirow{2}{\textbf{Fitur}} &
    \multicolumn{1}{c}{\rotatebox{90}{\textbf{SUPER ADMIN}}} &
    \multicolumn{1}{c}{\rotatebox{90}{\textbf{ADMIN}}} &
    \multicolumn{1}{c}{\rotatebox{90}{\textbf{PEMINJAM}}} \\
  \cmidrule{2-4}
  & \centering \textbf{\roleSuperAdmin} &
    \centering \textbf{\roleAdmin} &
    \centering \textbf{\rolePeminjam} \\
  \midrule
  \endhead

  \bottomrule
  \endfoot

  % ── Manajemen Pengguna ──
  \rowcolor{lightblue}\multicolumn{4}{l}{\textbf{A. Manajemen Pengguna}} \\[4pt]
  Mengelola akun Super Admin       & \cmark & \textendash & \textendash \\[4pt]
  Mengelola akun Admin             & \cmark & \textendash & \textendash \\[4pt]
  Mengelola akun Peminjam          & \cmark & \cmark      & \textendash \\[4pt]
  Menetapkan akses Satker admin   & \cmark & \textendash & \textendash \\[4pt]
  Mengubah informasi profil sendiri& \cmark & \cmark      & \cmark \\[4pt]
  Mengubah kata sandi sendiri      & \textendash$^\dagger$ & \cmark & \cmark \\[4pt]

  % ── Manajemen Barang ──
  \rowcolor{lightblue}\multicolumn{4}{l}{\textbf{B. Manajemen Barang}} \\[4pt]
  Menambah barang                  & \cmark & \cmark      & \textendash \\[4pt]
  Mengedit barang                  & \cmark & \cmark      & \textendash \\[4pt]
  Menghapus barang                 & \cmark & \cmark      & \textendash \\[4pt]
  Impor barang dari Excel          & \cmark & \cmark      & \textendash \\[4pt]
  Ekspor barang ke Excel           & \cmark & \cmark      & \textendash \\[4pt]
  Melihat katalog barang           & \cmark & \cmark      & \cmark \\[4pt]

  % ── Manajemen Peminjaman ──
  \rowcolor{lightblue}\multicolumn{4}{l}{\textbf{C. Manajemen Peminjaman}} \\[4pt]
  Mengajukan peminjaman            & \cmark$^\ddagger$ & \cmark$^\ddagger$ & \cmark \\[4pt]
  Menyetujui peminjaman            & \cmark & \cmark      & \textendash \\[4pt]
  Menolak peminjaman               & \cmark & \cmark      & \textendash \\[4pt]
  Menandoveran barang              & \cmark & \cmark      & \textendash \\[4pt]
  Memproses pengembalian           & \cmark & \cmark      & \textendash \\[4pt]
  Membuat peminjaman atas nama     & \cmark & \cmark      & \textendash \\[4pt]
  Melihat seluruh peminjaman       & \cmark & \cmark      & \textendash \\[4pt]
  Melihat riwayat pribadi          & \cmark & \cmark      & \cmark \\[4pt]
  Memindai QR Code                 & \cmark & \cmark      & \textendash \\[4pt]

  % ── Manajemen Satker ──
  \rowcolor{lightblue}\multicolumn{4}{l}{\textbf{D. Manajemen Satker}} \\[4pt]
  CRUD Satker                      & \cmark & \textendash & \textendash \\[4pt]

  % ── Log & Audit ──
  \rowcolor{lightblue}\multicolumn{4}{l}{\textbf{E. Log \& Audit}} \\[4pt]
  Melihat log audit                & \cmark & \cmark      & \textendash \\[4pt]
  Melihat riwayat impor             & \cmark & \cmark      & \textendash \\[4pt]

  % ── Notifikasi ──
  \rowcolor{lightblue}\multicolumn{4}{l}{\textbf{F. Notifikasi}} \\[4pt]
  Menerima notifikasi              & \cmark & \cmark      & \cmark \\[4pt]
  Mengelola notifikasi             & \cmark & \cmark      & \cmark \\[4pt]

  % ── Lingkup Akses ──
  \rowcolor{lightblue}\multicolumn{4}{l}{\textbf{G. Lingkup Akses}} \\[4pt]
  Akses lintas Satker              & \cmark & \textendash & \textendash \\[4pt]
  Akses terbatas pada Satker tugas& \textendash & \cmark & \textendash \\[4pt]
  Akses data pribadi saja          & \textendash & \textendash & \cmark \\[4pt]
\end{longtable}

\vspace{6pt}
\small
\noindent\textbf{Keterangan}:
\begin{itemize}
  \item[\cmark] \ Akses penuh
  \item[\textendash] \ Tidak memiliki akses
  \item[$^\dagger$] Super Admin dikecualikan dari operasi perubahan
        kata sandi diri sendiri karena alasan keamanan.
  \item[$^\ddagger$] Super Admin dan Admin yang memiliki peran
        {\rolePeminjam} juga dapat mengajukan peminjaman.
\end{itemize}
\normalsize

\section{Cakupan Akses Berdasarkan Satker}
\index{satker scoping}
\index{satuan kerja}
\index{satker access}
Sistem menggunakan konsep {\it Satker} (Satuan Kerja) sebagai unit
pembatasan akses data:

\begin{itemize}
  \item {\roleSuperAdmin} memiliki akses tanpa batasan Satker ({\it
        cross-Satker}). Semua data dari seluruh Satker dapat dilihat
        dan dikelola.
  \item {\roleAdmin} hanya dapat mengakses dan mengelola data yang
        terkait dengan Satker yang tercantum dalam kolom
        \texttt{satkerAkses} pada akunnya. Data Satker lain tidak
        dapat dilihat.
  \item {\rolePeminjam} hanya dapat melihat dan berinteraksi dengan
        data barang dan transaksi yang berasal dari Satker yang sama
        dengan profilnya.
\end{itemize}

% ─── Chapter 4: Fitur Super Admin ───────────────────────────────────────────
\chapter{Fitur dan Halaman --- Super Admin}

\index{super admin}
\index{halaman super admin}
\verb|/super-admin| --- Lingkup akses: seluruh Satker (tanpa batasan).

\section{Dashboard}
\index{dashboard super admin}
\textbf{Rute}: \texttt{/super-admin/dashboard}

Halaman utama yang menyajikan statistik lintas Satker, meliputi:
\begin{itemize}
  \item Total jumlah barang keseluruhan;
  \item Rincian stok barang berdasarkan kondisi (Baik, Rusak Ringan,
        Rusak Berat);
  \item Jumlah peminjaman berdasarkan status (Menunggu, Disetujui,
        Aktif/Dipinjam, Terlambat, Selesai);
  \item Perincian statistik per-Satker; dan
  \item Visualisasi data dalam bentuk grafik dan diagram.
\end{itemize}

\section{Manajemen Admin}
\index{manajemen admin}
\textbf{Rute}: \texttt{/super-admin/admin}

Halaman untuk mengelola seluruh akun Administrator, dengan kemampuan:
\begin{itemize}
  \item {\bfseries Membuat} akun admin baru;
  \item {\bfseries Mengaktifkan/Menonaktifkan} akun;
  \item {\bfseries Mengubah peran} (menaikkan ke Super Admin atau
        menurunkan ke Peminjam);
  \item {\bfseries Menetapkan akses Satker} untuk setiap admin,
        sehingga admin hanya dapat mengelola data di Satker
        yang telah ditentukan; dan
  \item {\bfseries Mereset kata sandi} akun admin.
\end{itemize}

\section{Katalog Barang (Lintas Satker)}
\index{katalog barang super admin}
\textbf{Rute}: \texttt{/super-admin/barang}

Menampilkan seluruh data barang BMN dari semua Satker. Super Admin
dapat melihat, mencari, dan menyaring barang berdasarkan:
kode Satker, kategori barang, kondisi, lokasi penyimpanan,
dan kata kunci pencarian. Navigasi detail barang tersedia di
\texttt{/super-admin/barang/[id]}.

\section{Manajemen Pengguna (Peminjam)}
\index{manajemen pengguna}
\index{peminjam}
\textbf{Rute}: \texttt{/super-admin/pengguna}

Mengelola seluruh akun Peminjam di seluruh Satker, meliputi:
\begin{itemize}
  \item Melihat daftar pengguna dan informasi profil;
  \item Menonaktifkan atau mengaktifkan akun;
  \item Mereset kata sandi; dan
  \item Mengelola peran pengguna.
\end{itemize}

\section{Daftar Peminjaman (Lintas Satker)}
\index{daftar peminjaman}
\textbf{Rute}: \texttt{/super-admin/peminjaman}

Tabel lengkap seluruh transaksi peminjaman dari semua Satker
dengan fitur:
\begin{itemize}
  \item Penyaringan ({\it filter}) berdasarkan status peminjaman;
  \item Pencarian berdasarkan kode transaksi, nama peminjam,
        atau nama barang;
  \item Navigasi ke halaman detail masing-masing transaksi; dan
  \item Tindakan massal ({\it bulk actions}) seperti persetujuan
      secara massal, penandoveran secara massal, dan pengembalian secara massal.
\end{itemize}

Detail transaksi dapat dilihat di
\texttt{/super-admin/peminjaman/[id]}.

\section{Manajemen Satker}
\index{manajemen satker}
\index{satuan kerja}
\textbf{Rute}: \texttt{/super-admin/satker}

Laman untuk melakukan operasi CRUD ({\it Create, Read, Update, Delete})
terhadap entitas Satker (Satuan Kerja), meliputi:
\begin{itemize}
  \item {\bfseries Membuat} Satker baru dengan mengisi kode unik,
        nama lengkap, nama singkat, dan status aktif;
  \item {\bfseries Mengubah} informasi Satker;
  \item {\bfseries Mengaktifkan/Menonaktifkan} Satker; dan
  \item {\bfseries Menghapus} Satker yang tidak diperlukan
        (dengan validasi tidak ada data terkait).
\end{itemize}

\section{Log Audit}
\index{log audit}
\index{audit trail}
\textbf{Rute}: \texttt{/super-admin/logs}

Halaman untuk melihat seluruh riwayat aktivitas pengguna dalam sistem,
yang mencakup:
\begin{itemize}
  \item {\bfseries Aktivitas yang dicatat}: login/logout, pembuatan
        data, perubahan data (dengan snapshot data lama dan baru),
        penghapusan data, perubahan status peminjaman, impor data,
        dan lain-lain;
  \item {\bfseries Informasi chaque log}: timestamp kejadian, alamat
        IP sumber, {\it user agent} browser, dan pengguna yang
        melakukan tindakan;
  \item {\bfseries Penyaringan}: berdasarkan tipe aktivitas, entitas
        yang terlibat, rentang tanggal, dan pengguna tertentu; dan
  \item {\bfseries Statistik}: ringkasan jumlah aktivitas per kategori.
\end{itemize}

Data log audit disimpan dengan pendekatan {\it snapshot-free join},
yaitu seluruh informasi yang diperlukan dicatat langsung pada saat
kejadian ({\it denormalized}) untuk menjamin akurasi historis
meskipun data utama berubah atau dihapus.

% ─── Chapter 5: Fitur Admin ──────────────────────────────────────────────────
\chapter{Fitur dan Halaman --- Administrator}

\index{admin}
\index{halaman admin}
\index{manajemen barang}
\verb|/admin| --- Lingkup akses: terbatas pada Satker yang diberikan.

\section{Dashboard}
\index{dashboard admin}
\textbf{Rute}: \texttt{/admin/dashboard}

Halaman utama Administrator yang menampilkan ringkasan statistik
berbasis Satker yang dikelolanya:
\begin{itemize}
  \item Jumlah total barang dan rincian stok;
  \item Jumlah pengajuan yang menunggu persetujuan;
  \item Jumlah peminjaman aktif;
  \item Jumlah barang yang terlambat dikembalikan;
  \item Lima transaksi terakhir; dan
  \item Visualisasi data dalam bentuk grafik.
\end{itemize}

Navigasi ke detail kategori spesifik tersedia di
\texttt{/admin/dashboard/kategori/[kategori]}, dengan opsi
kategori: {\it barang}, {\it pengajuan\_menunggu},
{\it peminjaman\_aktif}, {\it barang\_terlambat},
dan {\it peminjam}.

\section{Katalog Barang}
\index{katalog barang admin}
\textbf{Rute}: \texttt{/admin/barang}

Menampilkan daftar barang BMN yang dimiliki oleh Satker Administrator.
Fitur pencarian dan penyaringan tersedia berdasarkan kategori,
kondisi, lokasi penyimpanan, dan kata kunci.

\subsection{Tambah Barang Tunggal}
\index{tambah barang}
\textbf{Rute}: \texttt{/admin/barang/tambah}

Formulir untuk menambahkan satu item barang baru secara manual
dengan mengisi seluruh field yang diperlukan.

\subsection{Tambah Barang Massal (Bulk)}
\index{tambah barang massal}
\index{bulk insert}
\textbf{Rute}: \texttt{/admin/barang/bulk}

Formulir untuk menambahkan beberapa barang sekaligus secara massal.
Sistem secara otomatis menghasilkan nomor NUP (Nomor Urut Barang)
yang berurutan sesuai dengan jumlah yang ditentukan oleh pengguna.

\subsection{Detail Barang}
\index{detail barang}
\textbf{Rute}: \texttt{/admin/barang/[id]}

Halaman informasi lengkap satu item barang yang juga menyediakan
fungsi {\it Edit} untuk memperbarui data barang.

\section{Daftar Peminjaman}
\index{daftar peminjaman admin}
\index{persetujuan}
\index{pengembalian}
\textbf{Rute}: \texttt{/admin/peminjaman}

Tabel seluruh transaksi peminjaman di lingkup Satker Administrator
dengan fitur {\it filter} berdasarkan status. Tindakan yang tersedia
meliputi:
\begin{itemize}
  \item {\bfseries Persetujuan} pengajuan (\texttt{MENUNGGU} $\rightarrow$
        \texttt{DISETUJUI});
  \item {\bfseries Penolakan} pengajuan (\texttt{MENUNGGU} $\rightarrow$
        \texttt{DITOLAK}) dengan kolom catatan;
  \item {\bfseries Penandoveran} barang secara fisik (\texttt{DISETUJUI}
        $\rightarrow$ \texttt{DIPINJAM});
  \item {\bfseries Konfirmasi pengembalian} (\texttt{DIPINJAM}/
        \texttt{TERLAMBAT} $\rightarrow$ \texttt{DIKEMBALIKAN});
  \item {\bfseries Pembuatan peminjaman} atas nama borrower tertentu
        (\texttt{/admin/peminjaman/buat}); dan
  \item {\bfseries Tindakan massal} ({\it bulk approve, handover, return,
        delete}) untuk efisiensi operasional.
\end{itemize}

Detail transaksi dapat dilihat di
\texttt{/admin/peminjaman/[id]}.

\section{Pemindaian QR Code}
\index{qr code scanner}
\index{scan qr}
\textbf{Rute}: \texttt{/admin/scan}

Halaman khusus untuk Administrator memindai kode QR pada dokumen
Surat Pernyataan Peminjaman borrower. Pemindaian dapat dilakukan
melalui:
\begin{itemize}
  \item {\bfseries Kamera} perangkat (pemindaian {\it real-time}
        menggunakan pustaka {\it html5-qrcode}); atau
  \item {\bfseries Masukan manual} kode QR jika kamera tidak tersedia.
\end{itemize}

Hasil pemindaian akan langsung menavigasi Administrator ke halaman
detail transaksi terkait untuk mempercepat proses verifikasi dan
pengembalian.

\section{Riwayat Impor}
\index{riwayat impor}
\index{import log}
\textbf{Rute}: \texttt{/admin/import-log}

Halaman yang menampilkan daftar seluruh riwayat operasi impor data
Excel yang pernah dilakukan oleh Administrator di Satkernya,
disertai statistik rinci yang mencakup:
\begin{itemize}
  \item Jumlah data yang berhasil ditambahkan;
  \item Jumlah data yang diperbarui;
  \item Jumlah data yang dilewati ({\it skipped}); dan
  \item Jumlah data yang gagal beserta alasan kegagalan per item.
\end{itemize}

% ─── Chapter 6: Fitur Peminjam ──────────────────────────────────────────────
\chapter{Fitur dan Halaman --- Peminjam}

\index{peminjam}
\index{halaman peminjam}
\index{borrower}
\verb|/peminjam| --- Lingkup akses: data pribadi dan transaksi sendiri.

\section{Dashboard}
\index{dashboard peminjam}
\textbf{Rute}: \texttt{/peminjam/dashboard}

Halaman utama Borrower yang menyajikan ringkasan pribadi:
\begin{itemize}
  \item Total jumlah pengajuan peminjaman;
  \item Jumlah peminjaman aktif ({\it on loan});
  \item Jumlah peminjaman yang telah selesai;
  \item Daftar peminjaman aktif terbaru; dan
  \item Informasi penting terkait batas maksimal peminjaman aktif.
\end{itemize}

\section{Katalog Barang}
\index{katalog barang peminjam}
\index{browse catalog}
\textbf{Rute}: \texttt{/peminjam/katalog}

Halaman penjelajahan ({\it browse}) seluruh barang BMN yang tersedia
untuk Satker borrower, dengan fitur:
\begin{itemize}
  \item Pencarian berdasarkan nama, merk, atau tipe barang;
  \item Penyaringan berdasarkan kategori barang; dan
  \item Navigasi ke halaman detail barang.
\end{itemize}

Detail barang dapat dilihat di \texttt{/peminjam/katalog/[id]},
termasuk informasi ketersediaan stok terkini.

\section{Keranjang}
\index{keranjang}
\index{cart}
\index{shopping cart}
\textbf{Rute}: \texttt{/peminjam/keranjang}

Sistem keranjang belanja ({\it shopping cart}) yang memungkinkan
borrower memilih beberapa barang sebelum mengajukan peminjaman.
Fitur utama meliputi:
\begin{itemize}
  \item {\bfseries Penambahan/penghapusan} barang dari keranjang;
  \item {\bfseries Penyesuaian jumlah} barang yang ingin dipinjam;
  \item {\bfseries {\it Real-time stock polling}}: sistem secara berkala
        memperbarui informasi stok untuk mencegah pemesanan ganda
        ({\it double-booking}) selama sesi keranjang berlangsung; dan
  \item {\bfseries Validasi otomatis} --- barang yang stoknya tidak
        tersedia akan ditandai atau dihapus secara otomatis dari
        keranjang.
\end{itemize}

\section{Pengajuan Peminjaman}
\index{pengajuan peminjaman}
\index{submit borrowing request}
\index{draft mode}
\index{mode konsep}
\textbf{Rute}: \texttt{/peminjam/ajukan}

Formulir utama untuk mengajukan permintaan peminjaman. Proses ini
mendukung dua mode:

\begin{enumerate}
  \item {\bfseries Mode Konsep ({\it Draft Mode})}: Borrower dapat
        menyimpan pengajuan {\it tanpa} mengunggah Surat Pernyataan
        yang sudah ditandatangani. Item barang dikunci
        (tidak bisa dipinjam oleh pengguna lain) namun belum ada
        pengurangan stok. Pengajuan konsep dapat dilanjutkan dan
        diselesaikan kapan saja sebelum batas waktu berlaku.

  \item {\bfseries Mode Lengkap}: Borrower melengkapi seluruh
        informasi pengajuan termasuk:
        \begin{itemize}
          \item tanggal rencana pinjam dan tanggal rencana kembali;
          \item alasan/alasan peminjaman;
          \item pangkat dan golongan; dan
          \item {\it upload} Surat Pernyataan yang telah dicetak,
                ditandatangani, distempel, dan diunggah
                ({\it scan}/\it{foto} PDF).
        \end{itemize}
        Pada mode ini, status pengajuan berubah menjadi
        \texttt{MENUNGGU} dan notifikasi dikirimkan kepada seluruh
        Administrator terkait.
\end{enumerate}

Sistem membatasi setiap borrower maksimal 3 (tiga) peminjaman aktif
({\it maximum 3 active borrowings}) pada satu waktu. Batas ini
dapat dikonfigurasi melalui variabel lingkungan
\texttt{MAX\_PEMINJAMAN\_AKTIF}.

\section{Riwayat Peminjaman}
\index{riwayat peminjaman}
\index{borrowing history}
\textbf{Rute}: \texttt{/peminjam/riwayat}

Daftar seluruh riwayat pengajuan dan peminjaman borrower secara
pribadi, dengan fitur {\it filter} berdasarkan status:
{\it Menunggu}, {\it Disetujui}, {\it Dipinjam}, {\it Terlambat},
dan {\it Selesai}.

Detail setiap transaksi dapat dilihat di
\texttt{/peminjam/riwayat/[id]}, yang menyediakan:
\begin{itemize}
  \item {\bfseries Timeline} lengkap status pengajuan;
  \item {\bfseries Unduhan Surat Pernyataan} Peminjaman (PDF);
  \item {\bfseries Unduhan Surat Pengembalian} (jika sudah dikembalikan);
  \item {\bfseries Form ulasan/review} barang setelah pengembalian
        (\texttt{/peminjam/riwayat/[id]/review}).
\end{itemize}

% ─── Chapter 7: Fitur Bersama ───────────────────────────────────────────────
\chapter{Halaman dan Fitur Bersama}

\index{halaman bersama}
\index{shared pages}
\index{autentikasi}
Seluruh halaman di bawah ini dapat diakses oleh semua peran yang
telah {\it login} ke dalam sistem.

\section{Halaman Publik}
\index{halaman publik}

\subsection{Landing Page}
\index{landing page}
\textbf{Rute}: \texttt{/}

Halaman muka ({\it hero page}) yang menyambut pengunjung dengan
informasi singkat sistem dan tombol navigasi ke halaman login
maupun registrasi.

\subsection{Login}
\index{halaman login}
\textbf{Rute}: \texttt{/login}

Halaman autentikasi pengguna. Pengguna dengan {\it multi-role}
akan melihat pilihan peran sebelum masuk ke sistem. {\it Flow}
login meliputi:
\begin{enumerate}
  \item Masukan email dan kata sandi;
  \item {\it CSRF token} divalidasi;
  \item JWT Access Token (berlaku 60 menit) dan Refresh Token
        (via httpOnly cookie, berlaku 7 hari) diterbitkan;
  \item {\it Token versioning} diperbarui untuk menginvalidasi
        sesi sebelumnya secara otomatis; dan
  \item Pengalihan ke halaman sesuai peran.
\end{enumerate}

\subsection{Registrasi}
\index{halaman registrasi}
\textbf{Rute}: \texttt{/register}

Formulir pendaftaran untuk pengguna baru. Registrasi publik
secara {\it default} menghasilkan akun dengan peran
{\rolePeminjam} saja. Admin dan Super Admin harus dibuatkan
akunnya oleh administrator yang berwenang melalui halaman
manajemen pengguna masing-masing.

\subsection{Panduan}
\index{panduan}
\index{user guide}
\textbf{Rute}: \texttt{/panduan}

Halaman bantuan berisi panduan penggunaan sistem bagi pengguna
akhir (end-user), yang menjelaskan langkah-langkah pengajuan
peminjaman hingga pengembalian barang secara ringkas dan mudah
dipahami.

\section{Halaman Pengaturan}
\index{pengaturan}
\index{settings}
\index{profil}
\textbf{Rute}: \texttt{/pengaturan}

Halaman pengelolaan profil pengguna yang dapat diakses oleh
seluruh peran. Informasi yang dapat diubah meliputi:
\begin{itemize}
  \item Nama lengkap;
  \item Alamat email;
  \item NIP (Nomor Induk Pegawai);
  \item Jabatan;
  \item Unit Kerja;
  \item Eselon II, III, dan IV; dan
  \item Tipe laptop (apabila diperlukan).
\end{itemize}

Catatan: {\roleSuperAdmin} dikecualikan dari operasi perubahan
kata sandi sendiri karena alasan keamanan --- perubahan kata
sandi Super Admin harus dilakukan melalui mekanisme yang
terpisah dan lebih ketat.

\section{Halaman Notifikasi}
\index{notifikasi}
\index{notifications}
\textbf{Rute}: \texttt{/notifikasi}

Pusat pemberitahuan dalam aplikasi yang menyimpan seluruh notifikasi
pengguna. Setiap notifikasi memiliki:
\begin{itemize}
  \item {\bfseries Judul} dan {\bfseries Isi pesan};
  \item {\bfseries Prioritas}: Tinggi, Sedang, atau Rendah;
  \item {\bfseries Status}: Sudah dibaca atau belum;
  \item {\bfseries Tautan referensi} ke halaman terkait.
\end{itemize}

Pengguna dapat menandai notifikasi sebagai sudah dibaca,
menghapus notifikasi, dan mengakses halaman terkait dengan
satu klik.

\section{Halaman Bantuan}
\index{bantuan}
\index{help}
\textbf{Rute}: \texttt{/bantuan}

Halaman bantuan umum ({\it help center}) yang memuat:
\begin{itemize}
  \item Pertanyaan yang sering diajukan ({\it FAQ});
  \item Penjelasan alur kerja sistem; dan
  \item Informasi kontak teknis untuk pelaporan masalah.
\end{itemize}

% ─── Chapter 8: Alur Kerja Peminjaman ────────────────────────────────────────
\chapter{Alur Kerja Peminjaman}

\index{alur kerja}
\index{borrowing workflow}
\index{siklus hidup peminjaman}
\index{status peminjaman}

\section{Daftar Status Peminjaman}
\index{status peminjaman}

\begin{longtable}[!htbp]{p{3.5cm} p{4cm} p{6cm}}
  \caption{Daftar Status Peminjaman}
  \label{tab:status-peminjaman} \\

  \toprule
  \textbf{Status} & \textbf{Alias} & \textbf{Deskripsi} \\
  \midrule
  \endfirsthead

  \caption[]{Daftar Status Peminjaman (sambungan)} \\
  \toprule
  \textbf{Status} & \textbf{Alias} & \textbf{Deskripsi} \\
  \midrule
  \endhead

  \bottomrule
  \endfoot

  DRAFT           & Konsep      & Pengajuan disimpan tanpa surat
                   ditandatangani. Barang dikunci namun stok belum
                   dikurangi. Dapat dilanjutkan kapan saja. \\[4pt]
  MENUNGGU        & Menunggu    & Pengajuan lengkap telah disubmit.
                   Menunggu persetujuan Administrator. \\[4pt]
  DISETUJUI       & Disetujui   & Administrator menyetujui pengajuan.
                   Stok dikurangi, QR Code dibangkitkan. \\[4pt]
  DITOLAK         & Ditolak     & Administrator menolak pengajuan
                   dengan alasan tertentu. \\[4pt]
  DIPINJAM        & Dipinjam    & Barang secara fisik telah handed over
                   kepada borrower. \\[4pt]
  TERLAMBAT       & Terlambat   & Tanggal kembali rencana telah melewati
                   batas tanpa pengembalian. \\[4pt]
  DIKEMBALIKAN   & Selesai     & Barang telah dikembalikan dan
                   diverifikasi oleh Administrator. \\
\end{longtable}

\section{Diagram Alur Status}
\index{diagram alur}
\index{state diagram}

Diagram transisi status peminjaman digambarkan pada
Gambar~\ref{fig:status-flow}.

\begin{figure}[!htbp]
  \centering
  % Since we can't draw actual graphics, we'll describe the flow textually
  % and create a simple text-based representation
  \IfFileExists{./assets/status-flow.png}
    {\includegraphics[width=0.85\textwidth]{./assets/status-flow.png}}
    {%
      \begin{tabularx}{0.85\textwidth}{C C C}
        \toprule
        \multicolumn{3}{c}{\textbf{Draft}} \\[8pt]
        \downarrow{\small submit} &
        \multicolumn{2}{c}{} \\[4pt]
        \multicolumn{3}{c}{\textbf{Menunggu}} \\[8pt]
        \begin{tabular}{c}
          approve \\ \hline
          \downarrow \\ reject \\ \end{tabular} &
        \begin{tabular}{c}
          \rightarrow \\ reject \\ \end{tabular} \\[14pt]
        \textbf{Disetujui} &
        \textbf{Ditolak} &
        \\[8pt]
        handover \downarrow &
        \multicolumn{2}{c}{} \\[4pt]
        \multicolumn{2}{c}{\textbf{Dipinjam}} \\[8pt]
        return \downarrow &
        late \downarrow \\[4pt]
        \textbf{Dikembalikan} &
        \textbf{Terlambat} \\[8pt]
      \end{tabularx}
    }
  \caption{Diagram Alur Transisi Status Peminjaman}
  \label{fig:status-flow}
\end{figure}

\section{Penjelasan Tahapan Proses}
\index{tahapan proses}
\index{process stages}

\subsection{Tahap 1: Konsep (Draft)}
\index{draft mode}
Borrower menyusun pengajuan peminjaman di halaman
\texttt{/peminjam/ajukan}. Pada tahap ini:
\begin{itemize}
  \item Informasi barang, tanggal rencana, alasan, dan data
        lainnya dapat disimpan tanpa mengunggah Surat Pernyataan;
  \item Barang yang ditambahkan ke pengajuan {\it dikunci}
        ({\it soft lock}) agar tidak bisa diajukan oleh borrower
        lain. Namun, {\it stock count} belum dikurangi;
  \item Pengajuan konsep dapat diedit, dilengkapi, atau
        dibatalkan kapan saja.
\end{itemize}

\subsection{Tahap 2: Pengajuan (Menunggu)}
\index{menunggu persetujuan}
Setelah borrower melengkapi seluruh informasi dan mengunggah
Surat Pernyataan yang telah ditandatangani, status berubah
menjadi \texttt{MENUNGGU}. Sistem secara otomatis:
\begin{itemize}
  \item Mengirim {\bfseries notifikasi} kepada seluruh Administrator
        di Satker terkait;
  \item {\bfseries Mematikan tombol batal} untuk pengajuan
        pada sisi borrower (pengajuan tidak bisa dibatalkan
        lagi setelah disubmit).
\end{itemize}

\subsection{Tahap 3: Persetujuan (Disetujui / Ditolak)}
\index{persetujuan}
\index{penolakan}
Administrator di Satker terkait meninjau pengajuan. Terdapat
dua kemungkinan hasil:

\begin{itemize}
  \item {\bfseries Disetujui} (\texttt{MENUNGGU} $\rightarrow$
        \texttt{DISETUJUI}):
        \begin{itemize}
          \item Stok barang dikurangi secara atomik;
          \item {\it QR Code} dibangkitkan dalam warna biru tua
                (\#1e3a5f) pada latar putih dan diunggah ke Vercel Blob;
          \item Nomor Surat Pernyataan dibuat secara berurutan
                dengan format \texttt{PRN-<nomor>/BMN/PP.1/<tahun>};
          \item Notifikasi dikirimkan kepada borrower.
        \end{itemize}
  \item {\bfseries Ditolak} (\texttt{MENUNGGU} $\rightarrow$
        \texttt{DITOLAK}):
        \begin{itemize}
          \item Kunci pada barang dilepaskan (stok tidak berubah);
          \item Catatan penolakan dari Administrator dicatat;
          \item Notifikasi penolakan dikirimkan kepada borrower.
        \end{itemize}
\end{itemize}

\subsection{Tahap 4: Penandoveran (Dipinjam)}
\index{penandoveran}
\index{handover}
Borrower yang pengajuannya disetujui mengambil barang secara
fisik di lokasi penyimpanan. Administrator mengonfirmasi
penandoveran (\texttt{DISETUJUI} $\rightarrow$ \texttt{DIPINJAM}).

\subsection{Tahap 5: Pengembalian}
\index{pengembalian}
\index{return process}

Proses pengembalian terdiri dari dua langkah:

\begin{enumerate}
  \item {\bfseries Borrower meminta pengembalian}: Borrower
        mengunggah Surat Pengembalian yang telah ditandatangani
        melalui halaman riwayat (\texttt{/peminjam/riwayat/[id]}).
        Notifikasi dikirim kepada Administrator.

  \item {\bfseries Administrator mengonfirmasi pengembalian}:
        (\texttt{DIPINJAM}/\texttt{TERLAMBAT} $\rightarrow$
        \texttt{DIKEMBALIKAN}):
        \begin{itemize}
          \item Stok barang dipulihkan;
          \item {\it Catatan pengembalian} dicatat oleh Administrator;
          \item Surat Pengembalian dengan nomor urut, perangko
                digital, dan tanda tangan digital Administrator
                dibangkitkan; dan
          \item Borrower dapat memberikan {\it review} untuk
                barang yang dikembalikan
                (\texttt{/peminjam/riwayat/[id]/review}).
        \end{itemize}
\end{enumerate}

\section{Pencegahan Pemesanan Ganda}
\index{double-booking prevention}
\index{pemesanan ganda}
\index{stock lock}

Sistem menerapkan mekanisme {\it pessimistic locking} untuk
mencegah {\it double-booking}:

\begin{itemize}
  \item Barang yang berada dalam status \texttt{DRAFT},
        \texttt{MENUNGGU}, \texttt{DISETUJUI}, \texttt{DIPINJAM},
        atau \texttt{TERLAMBAT} tidak dapat disertakan dalam
        pengajuan baru;
  \item {\it Real-time stock polling} pada halaman keranjang
        memastikan borrower selalu melihat informasi stok terkini;
  \item Setiap operasi pengurangan dan pemulihan stok dilakukan
        secara {\it atomik} dalam transaksi basis data
        menggunakan Prisma ORM.
\end{itemize}

\section{Deteksi Terlambat Otomatis}
\index{detection otomatis}
\index{auto late detection}
\index{terlambat}
\index{overdue}

Sistem secara otomatis menandai peminjaman sebagai \texttt{TERLAMBAT}
ketika tanggal kembali rencana (\texttt{tanggalKembaliRencana})
telah dilewati. Deteksi ini dilakukan secara {\it throttled}
(maksimal sekali per menit per {\it dashboard load}) untuk
menghindari {\it performance overhead} pada basis data.

\section{Batas Maksimum Peminjaman Aktif}
\index{maksimum peminjaman aktif}
\index{max active borrowings}
\index{active loan limit}

Setiap borrower dibatasi maksimal 3 (tiga) peminjaman aktif
({\it active borrowings}) secara bersamaan. Batas ini
dapat dikonfigurasi melalui variabel lingkungan
\texttt{MAX\_PEMINJAMAN\_AKTIF}. Apabila batas telah tercapai,
borrower tidak dapat mengajukan peminjaman baru sebelum ada
yang ditandai selesai.

% ─── Chapter 9: Entitas Basis Data ─────────────────────────────────────────
\chapter{Entitas Basis Data}

\index{basis data}
\index{database entities}
\index{schema}
\index{prisma}

Sistem menggunakan Prisma ORM dengan basis data PostgreSQL.
Berikut adalah entitas utama yang menyusun skema basis data.

\section{ERD Singkat}
\index{ERD}
\index{entity relationship diagram}

Model hubungan antar-entitas utama digambarkan pada
Gambar~\ref{fig:erd}.

\begin{figure}[htbp]
  \centering
  \IfFileExists{./assets/erd.png}
    {\includegraphics[width=0.9\textwidth]{./assets/erd.png}}
    {%
      \begin{tabular}{|p{2.5cm}|p{2.5cm}|p{2.5cm}|p{2.5cm}|}
        \hline
        \textbf{Satker} & \textbf{User} & \textbf{Barang} & \textbf{Peminjaman} \\
        \hline
        kode (PK) & id (PK) & kodeBarang (PK) & id (PK) \\
        nama & nip & nama & kodeTransaksi \\
        singkat & email & merk & userId (FK) \\
        aktif & roles & jenis & status \\
        & satkerAkses & jumlahTotal & tanggalPengajuan \\
        & & jumlahTersedia & tanggalPinjamRencana \\
        & & kondisi & disetujuiOleh (FK) \\
        & & kodeSatker (FK) & dikembalikanOleh (FK) \\
        \hline
        \multicolumn{4}{|c|}{\textbf{DetailPeminjaman}} \\
        \hline
        \multicolumn{4}{|c|}{id, peminjamanId, barangKode, jumlahPinjam, statusItem} \\
        \hline
        \textbf{Notifikasi} & \textbf{AuditLog} & \textbf{ImportLog} & \textbf{NomorSuratCounter} \\
        \hline
        id, tipe, judul, pesan, prioritas, isBaca, referenceId &
        id, aksi, entitas, entitasId, dataLama, dataBaru, IP, userAgent, timestamp &
        id, jenis, statistik, adminId &
        jenis, tahun, counter \\
        \hline
      \end{tabular}
    }
  \caption{Hubungan Antar-Entitas Utama}
  \label{fig:erd}
\end{figure}

\section{Deskripsi Entitas}
\index{deskripsi entitas}

\subsection{Satker (Satuan Kerja)}
\index{entitas!Satker}
\index{satuan kerja}
\index{cost center}

\begin{longtable}[!htbp]{p{3.5cm} p{9cm}}
  \caption{Entitas Satker}
  \label{tab:entitas-satker} \\
  \toprule
  \textbf{Atribut} & \textbf{Deskripsi} \\
  \midrule
  \endfirsthead
  \caption[]{Entitas Satker (sambungan)} \\
  \toprule
  \textbf{Atribut} & \textbf{Deskripsi} \\
  \midrule
  \endhead
  \bottomrule
  \endfoot

  \texttt{kode}    & Kode unik Satker (Primary Key). \\
  \texttt{nama}    & Nama lengkap Satuan Kerja. \\
  \texttt{singkat}& Nama singkat/gelar Satker. \\
  \texttt{aktif}  & Status aktif Satker (boolean).\\
\end{longtable}

\subsection{User (Pengguna)}
\index{entitas!User}
\index{pengguna}
\index{user account}

\begin{longtable}[!htbp]{p{3.5cm} p{9cm}}
  \caption{Entitas User}
  \label{tab:entitas-user} \\
  \toprule
  \textbf{Atribut} & \textbf{Deskripsi} \\
  \midrule
  \endfirsthead
  \caption[]{Entitas User (sambungan)} \\
  \toprule
  \textbf{Atribut} & \textbf{Deskripsi} \\
  \midrule
  \endhead
  \bottomrule
  \endfoot

  \texttt{id}                & ID unik pengguna (Primary Key, auto-increment). \\
  \texttt{nama}             & Nama lengkap pengguna. \\
  \texttt{nip}              & Nomor Induk Pegawai. \\
  \texttt{email}            & Alamat email (unik). \\
  \texttt{password}         & Hash kata sandi (bcrypt). \\
  \texttt{jabatan}          & Jabatan pengguna. \\
  \texttt{unitKerja}        & Unit kerja pengguna. \\
  \texttt{eselon2/3/4}      & Pengelompokan eselon. \\
  \texttt{tipeLaptop}       & Tipe laptop (opsional). \\
  \texttt{roles}            & Array enum peran: \texttt{[SUPER\_ADMIN, ADMIN, PEMINJAM]}.
                              Satu akun dapat memiliki beberapa peran. \\
  \texttt{sumber}           & Asal data: \texttt{MANUAL} atau \texttt{IMPORT}. \\
  \texttt{tokenVersion}     & Nomor versi token untuk invalidasi sesi. \\
  \texttt{retirementDate}   & Tanggal pensiun (dihitung otomatis dari NIP). \\
  \texttt{satkerAkses}      & Array kode Satker yang boleh diakses oleh Admin.
                              Super Admin mengabaikan batasan ini. \\
  \texttt{lastActivityAt}   & Timestamp aktivitas terakhir. \\
  \texttt{sessionInvalidatedAt} & Timestamp invalidasi sesi terakhir.
                              Semua sesi dengan timestamp lebih lama
                              dari nilai ini dianggap tidak valid. \\
\end{longtable}

\subsection{Barang}
\index{entitas!Barang}
\index{barang}
\index{barang milik negara}
\index{BMN}

\begin{longtable}[!htbp]{p{3.5cm} p{9cm}}
  \caption{Entitas Barang}
  \label{tab:entitas-barang} \\
  \toprule
  \textbf{Atribut} & \textbf{Deskripsi} \\
  \midrule
  \endfirsthead
  \caption[]{Entitas Barang (sambungan)} \\
  \toprule
  \textbf{Atribut} & \textbf{Deskripsi} \\
  \midrule
  \endhead
  \bottomrule
  \endfoot

  \texttt{kodeBarang}       & Kode unik barang = \texttt{<kodeSatker>-<kodeBarangBMN>-<NUP>}.
                              Kunci unik komposit \texttt{(kodeSatker,
                              kodeBarangBmn, nup)}. \\
  \texttt{nama}             & Nama barang. \\
  \texttt{merk}             & Merk/type/model barang. \\
  \texttt{tipe}             & Tipe/sub-tipe barang (field baru). \\
  \texttt{jenis}            & Kategori: \texttt{ELEKTRONIK}, \texttt{FURNITUR},
                              \texttt{KENDARAAN}, \texttt{ATK}, \texttt{LAINNYA}. \\
  \texttt{jumlahTotal}      & Total jumlah unit barang. \\
  \texttt{jumlahTersedia}   & Jumlah unit yang sedang tersedia (belum dipinjam). \\
  \texttt{kondisi}          & Kondisi barang: \texttt{BAIK}, \texttt{RUSAK\_RINGAN},
                              \texttt{RUSAK\_BERAT}. \\
  \texttt{lokasiPenyimpanan}& Lokasi penyimpanan barang. \\
  \texttt{deskripsi}        & Deskripsi/uraian barang. \\
  \texttt{fotoUrl}          & URL foto barang di Vercel Blob. \\
  \texttt{sumber}           & Asal data: \texttt{MANUAL} atau \texttt{IMPORT}. \\
  \texttt{kodeSatker}       & Kode Satker pemilik barang (FK). \\
  \texttt{namaSatker}       & Nama Satker (snapshot denormalisasi). \\
  \texttt{kodeBarangBmn}    & Kode barang menurut daftar BMN. \\
  \texttt{nup}              & Nomor Urut barang (bagian dari Primary Key). \\
\end{longtable}

\subsection{Peminjaman}
\index{entitas!Peminjaman}
\index{peminjaman}
\index{loan transaction}

\begin{longtable}[!htbp]{p{3.5cm} p{9cm}}
  \caption{Entitas Peminjaman}
  \label{tab:entitas-peminjaman} \\
  \toprule
  \textbf{Atribut} & \textbf{Deskripsi} \\
  \midrule
  \endfirsthead
  \caption[]{Entitas Peminjaman (sambungan)} \\
  \toprule
  \textbf{Atribut} & \textbf{Deskripsi} \\
  \midrule
  \endhead
  \bottomrule
  \endfoot

  \texttt{id}                        & ID unik transaksi (Primary Key). \\
  \texttt{kodeTransaksi}             & Kode unik transaksi (unique). \\
  \texttt{kodePeminjaman}            & Snapshot kode barang pada saat pengajuan. \\
  \texttt{userId}                    & ID borrower/peminjam (FK). \\
  \texttt{tanggalPengajuan}          & Tanggal pengajuan dibuat. \\
  \texttt{tanggalKirim}             & Tanggal submit (upload surat). \\
  \texttt{tanggalPinjamRencana}     & Tanggal rencana pinjam dimulai. \\
  \texttt{tanggalKembaliRencana}    & Batas tanggal pengembalian. \\
  \texttt{tanggalKembaliAktual}     & Tanggal pengembalian sesungguhnya. \\
  \texttt{tanggalPermintaanKembali} & Tanggal borrower meminta pengembalian. \\
  \texttt{status}                   & Status saat ini. \\
  \texttt{alasanPeminjaman}         & Alasan/tujuan peminjaman. \\
  \texttt{pangkatGolongan}          & Pangkat dan golongan borrower. \\
  \texttt{dokumenUrl}               & URL Surat Pernyataan ( borrower). \\
  \texttt{dokumenStempelUrl}        & URL Surat dengan perangko digital. \\
  \texttt{dokumenPengembalianUrl}   & URL Surat Pengembalian. \\
  \texttt{qrCodeUrl}               & URL QR Code transaksi. \\
  \texttt{catatanAdmin}             & Catatan dari Administrator (saat penolakan/approval). \\
  \texttt{catatanPengembalian}      & Catatan saat pengembalian. \\
  \texttt{disetujuiOleh}            & ID Admin yang menyetujui (FK, nullable). \\
  \texttt{dikembalikanOleh}         & ID Admin yang memproses pengembalian (FK, nullable). \\
  \texttt{nomorSurat}               & Nomor Surat Pernyataan (\texttt{PRN-<no>/BMN/PP.1/<tahun>}). \\
  \texttt{tahunSurat}               & Tahun surat. \\
\end{longtable}

\subsection{DetailPeminjaman}
\index{entitas!DetailPeminjaman}
\index{detil peminjaman}
\index{junction table}

\begin{longtable}[!htbp]{p{3.5cm} p{9cm}}
  \caption{Entitas DetailPeminjaman}
  \label{tab:entitas-detail-peminjaman} \\
  \toprule
  \textbf{Atribut} & \textbf{Deskripsi} \\
  \midrule
  \endfirsthead
  \caption[]{Entitas DetailPeminjaman (sambungan)} \\
  \toprule
  \textbf{Atribut} & \textbf{Deskripsi} \\
  \midrule
  \endhead
  \bottomrule
  \endfoot

  \texttt{id}            & ID unik (Primary Key). \\
  \texttt{peminjamanId}  & ID transaksi peminjaman (FK, cascade delete). \\
  \texttt{barangKode}   & Kode barang yang dipinjam (snapshot saat pengajuan). \\
  \texttt{jumlahPinjam} & Jumlah unit yang dipinjam. \\
  \texttt{statusItem}   & Status per-item: \texttt{DIPINJAM} atau
                          \texttt{DIKEMBALIKAN} (untuk pengembalian sebagian). \\
\end{longtable}

\subsection{Entitas Pendukung}
\index{entitas pendukung}

\begin{longtable}[!htbp]{p{3.5cm} p{3.5cm} p{5.5cm}}
  \caption{Entitas Pendukung}
  \label{tab:entitas-pendukung} \\
  \toprule
  \textbf{Entitas} & \textbf{Jumlah} & \textbf{Fungsi} \\
  \midrule
  \endfirsthead
  \caption[]{Entitas Pendukung (sambungan)} \\
  \toprule
  \textbf{Entitas} & \textbf{Jumlah} & \textbf{Fungsi} \\
  \midrule
  \endhead
  \bottomrule
  \endfoot

  NomoroSuratCounter & 1 & Counter atomik untuk nomor surat.
  Menjamin keunikan dan urutan nomor surat per jenis per tahun. \\[6pt]
  Notifikasi         & N & Pemberitahuan dalam aplikasi
  per pengguna dengan tipe, prioritas, dan tautan referensi. \\[6pt]
  AuditLog           & N & Pencatatan aktivitas: aksi, entitas,
  entitasId, dataLama, dataBaru, IP, user agent, timestamp. \\[6pt]
  ImportLog          & N & Riwayat operasi impor Excel
  dengan statistik per-item. \\[6pt]
  BlacklistedToken   & N & Refresh token yang telah
  diinvalidasi dengan timestamp kadaluarsa. \\
\end{longtable}

% ─── Chapter 10: API Routes ─────────────────────────────────────────────────
\chapter{Rute API}

\index{API}
\index{rest api}
\index{endpoint}
\index{api routes}

\section{Daftar Endpoint API}
\index{daftar endpoint}

Seluruh endpoint API berada di bawah prefix \texttt{/api}.
Berikut adalah ringkasan endpoint utama yang dikelompokkan
berdasarkan fungsinya.

\begin{longtable}[!htbp]{p{4cm} p{5cm} p{3cm} p{2.5cm}}
  \caption{Daftar Endpoint API}
  \label{tab:api-routes} \\

  \toprule
  \textbf{Grup} & \textbf{Rute} & \textbf{Metode} & \textbf{Akses} \\
  \midrule
  \endfirsthead

  \caption[]{Daftar Endpoint API (sambungan)} \\
  \toprule
  \textbf{Grup} & \textbf{Rute} & \textbf{Metode} & \textbf{Akses} \\
  \midrule
  \endhead

  \bottomrule
  \endfoot

  % ── Auth ──
  \rowcolor{lightblue}\multicolumn{4}{l}{\textbf{Autentikasi}} \\[4pt]
                   & \texttt{/api/auth/login}         & POST   & Publik \\
                   & \texttt{/api/auth/register}     & POST   & Publik \\
                   & \texttt{/api/auth/logout}       & POST   & Auth \\
                   & \texttt{/api/auth/refresh}      & POST   & Auth \\
                   & \texttt{/api/auth/switch-role}   & POST   & Auth \\
                   & \texttt{/api/auth/me}           & GET    & Auth \\
                   & \texttt{/api/auth/password}      & PUT    & Auth \\
                   & \texttt{/api/auth/csrf-token}   & GET    & Publik \\

  % ── Barang ──
  \rowcolor{lightblue}\multicolumn{4}{l}{\textbf{Barang}} \\[4pt]
                   & \texttt{/api/barang}             & GET    & Auth \\
                   & \texttt{/api/barang}             & POST   & Admin+ \\
                   & \texttt{/api/barang/[id]}        & GET    & Auth \\
                   & \texttt{/api/barang/[id]}        & PUT    & Admin+ \\
                   & \texttt{/api/barang/[id]}        & DELETE & Admin+ \\
                   & \texttt{/api/barang/import}      & POST   & Admin+ \\
                   & \texttt{/api/barang/template}    & GET    & Admin+ \\
                   & \texttt{/api/barang/stock}      & POST   & Auth \\
                   & \texttt{/api/barang/merk}       & GET    & Admin+ \\
                   & \texttt{/api/barang/nup-preview} & GET    & Admin+ \\

  % ── Peminjaman ──
  \rowcolor{lightblue}\multicolumn{4}{l}{\textbf{Peminjaman}} \\[4pt]
                   & \texttt{/api/peminjaman}        & GET    & Admin+ \\
                   & \texttt{/api/peminjaman}         & POST   & Auth \\
                   & \texttt{/api/peminjaman/bulk}    & POST   & Admin+ \\
                   & \texttt{/api/peminjaman/[id]}    & GET    & Auth \\
                   & \texttt{/api/peminjaman/[id]}    & PUT    & Auth \\
                   & \texttt{/api/peminjaman/[id]}    & DELETE & Admin+ \\
                   & \texttt{/api/peminjaman/[id]/approve}    & POST & Admin+ \\
                   & \texttt{/api/peminjaman/[id]/reject}     & POST & Admin+ \\
                   & \texttt{/api/peminjaman/[id]/handover}   & POST & Admin+ \\
                   & \texttt{/api/peminjaman/[id]/return}     & POST & Admin+ \\
                   & \texttt{/api/peminjaman/[id]/pdf}        & GET  & Auth \\
                   & \texttt{/api/peminjaman/[id]/stempel}    & POST & Admin+ \\
                   & \texttt{/api/peminjaman/[id]/request-return} & POST & Auth \\
                   & \texttt{/api/peminjaman/scan}             & POST & Admin+ \\

  % ── Dashboard ──
  \rowcolor{lightblue}\multicolumn{4}{l}{\textbf{Dashboard}} \\[4pt]
                   & \texttt{/api/dashboard}             & GET & Admin+ \\
                   & \texttt{/api/dashboard/user}        & GET & Auth \\
                   & \texttt{/api/dashboard/super}       & GET & Super \\
                   & \texttt{/api/dashboard/kategori}    & GET & Admin+ \\

  % ── Users ──
  \rowcolor{lightblue}\multicolumn{4}{l}{\textbf{Manajemen User}} \\[4pt]
                   & \texttt{/api/users}            & GET    & Admin+ \\
                   & \texttt{/api/users}            & POST   & Admin+ \\
                   & \texttt{/api/users/[id]}      & GET    & Admin+ \\
                   & \texttt{/api/users/[id]}      & PUT    & Admin+ \\
                   & \texttt{/api/users/[id]}      & DELETE & Admin+ \\
                   & \texttt{/api/users/[id]/reset-password} & POST & Admin+ \\
                   & \texttt{/api/users/promote}   & POST   & Super \\

  % ── Satker ──
  \rowcolor{lightblue}\multicolumn{4}{l}{\textbf{Satker}} \\[4pt]
                   & \texttt{/api/satker}           & GET    & Admin+ \\
                   & \texttt{/api/satker}           & POST   & Super \\
                   & \texttt{/api/satker/[id]}      & PUT    & Super \\
                   & \texttt{/api/satker/[id]}      & DELETE & Super \\

  % ── Notifications ──
  \rowcolor{lightblue}\multicolumn{4}{l}{\textbf{Notifikasi}} \\[4pt]
                   & \texttt{/api/notifikasi}       & GET    & Auth \\
                   & \texttt{/api/notifikasi/[id]}  & PUT    & Auth \\
                   & \texttt{/api/notifikasi/[id]}  & DELETE & Auth \\

  % ── Audit Logs ──
  \rowcolor{lightblue}\multicolumn{4}{l}{\textbf{Log Audit}} \\[4pt]
                   & \texttt{/api/audit-logs}       & GET    & Admin+ \\

  % ── Export ──
  \rowcolor{lightblue}\multicolumn{4}{l}{\textbf{Ekspor}} \\[4pt]
                   & \texttt{/api/export/barang}    & GET    & Admin+ \\
                   & \texttt{/api/export/peminjaman}& GET    & Admin+ \\
                   & \texttt{/api/export/users}     & GET    & Admin+ \\

  % ── Import ──
  \rowcolor{lightblue}\multicolumn{4}{l}{\textbf{Impor}} \\[4pt]
                   & \texttt{/api/import-log}      & GET    & Admin+ \\
                   & \texttt{/api/import-pegawai}  & POST   & Admin+ \\
                   & \texttt{/api/import-peminjam} & POST  & Admin+ \\
\end{longtable}

\vspace{6pt}
\small
\noindent\textbf{Keterangan kolom \textit{Akses}}:
\begin{itemize}
  \item[\texttt{Publik}] --- Tidak memerlukan autentikasi.
  \item[\texttt{Auth}] --- Memerlukan Access Token JWT yang valid.
  \item[\texttt{Admin+}] --- Memerlukan peran {\roleAdmin} atau
        {\roleSuperAdmin}.
  \item[\texttt{Super}] --- Memerlukan peran {\roleSuperAdmin}.
\end{itemize}
\normalsize

% ─── Chapter 11: Fitur Keamanan ────────────────────────────────────────────
\chapter{Fitur Keamanan}

\index{keamanan}
\index{security}
\index{autentikasi}
\index{otorisasi}

\section{Mekanisme Autentikasi}
\index{mekanisme autentikasi}
\index{JWT}
\index{token}

Sistem menerapkan {\it multi-layered authentication}:

\begin{itemize}
  \item {\bfseries JWT Access Token}: Berlaku 60 menit, disimpan
        di memori aplikasi (state management). Berfungsi sebagai
        {\it bearer token} pada setiap request API.

  \item {\bfseries Refresh Token}: Berlaku 7 hari, disimpan dalam
        cookie httpOnly yang tidak dapat diakses oleh JavaScript
        ({\it XSS protection}).

  \item {\bfseries CSRF Protection}: Setiap request state-changing
        ({\it POST, PUT, DELETE}) divalidasi menggunakan pola
        {\it Double Submit Cookie}.

  \item {\bfseries Token Versioning}: Setiap akun memiliki
        \texttt{tokenVersion} yang diperbarui setiap kali
        kata sandi berubah atau sesi diinvalidasi. Token lama
        secara otomatis menjadi tidak valid.

  \item {\bfseries Token Blacklisting}: Refresh token yang di-revoke
        dicatat dalam tabel \texttt{BlacklistedToken} dengan
        timestamp kadaluarsa. Setiap {\it refresh} divalidasi
        terhadap daftar hitam ini.

  \item {\bfseries Session Invalidation}: Kolom
        \texttt{sessionInvalidatedAt} pada tabel {\it User}
        memungkinkan invalidasi masif seluruh sesi milik satu
        pengguna secara instan (berguna saat terjadi insiden
        keamanan).
\end{itemize}

\section{Mekanisme Otorisasi}
\index{mekanisme otorisasi}
\index{RBAC}
\index{middleware}

Semua rute API dilindungi oleh {\it middleware otorisasi} yang:
\begin{enumerate}
  \item Mengekstrak dan memvalidasi Access Token dari {\it header}
        Authorization;
  \item Memverifikasi peran ({\it role}) pengguna terhadap
        requirement rute yang diminta; dan
  \item Menerapkan {\it scope checking} berbasis Satker untuk
        role {\roleAdmin}.
\end{enumerate}

\section{Proteksi Tambahan}
\index{proteksi keamanan}
\index{rate limiting}
\index{helmet}
\index{file upload}

\begin{itemize}
  \item {\bfseries Rate Limiting}: Batasan jumlah request per IP
        menggunakan {\it express-rate-limit} untuk mencegah
        serangan brute-force dan DoS.
  \item {\bfseries Security Headers}: {\it helmet} digunakan
        untuk menetapkan {\it HTTP security headers} yang
        direkomendasikan (HSTS, X-Frame-Options, CSP, dll.).
  \item {\bfseries Validasi Input}: Seluruh input divalidasi
        menggunakan skema Zod pada level {\it controller},
        sebelum diproses lebih lanjut.
  \item {\bfseries Validasi File Upload}: {\it Multer}
        membatasi tipe file dan ukuran upload untuk mencegah
        {\it malicious file upload}.
\end{itemize}

% ─── Chapter 12: Fitur Tambahan ────────────────────────────────────────────
\chapter{Fitur Tambahan}

\index{fitur tambahan}

\section{Generasi Dokumen Otomatis}
\index{generasi dokumen}
\index{pdf generation}
\index{surat pernyataan}
\index{surat pengembalian}

\subsection{Surat Pernyataan Peminjaman}
\index{surat pernyataan}
\index{borrowing statement}
\index{PRN}
Dihasilkan secara otomatis saat Administrator {\it approve}
pengajuan. Dokumen mencakup:
\begin{itemize}
  \item {\it Header} resmi/leterkop instansi;
  \item {\bfseries Nomor surat} dengan format
        \texttt{PRN-<nomor>/BMN/PP.1/<tahun>};
  \item {\bfseries Tabel barang} yang dipinjam;
  \item {\bfseries Pernyataan kewajiban} borrower (6 butir);
  \item {\bfseries Blok tanda tangan} borrower; dan
  \item {\bfseries QR Code} transaksi.
\end{itemize}

Sistem menyediakan {\it preview} PDF sebelum borrower
mengunggah hasil scan/tanda tangan.

\subsection{Surat Pengembalian}
\index{surat pengembalian}
\index{return letter}
\index{return statement}
Dihasilkan saat Administrator mengonfirmasi pengembalian.
Dokumen mencakup:
\begin{itemize}
  \item {\it Header} resmi/leterkop instansi;
  \item {\bfseries Nomor surat pengembalian};
  \item {\bfseries Daftar barang} yang dikembalikan beserta kondisi;
  \item {\bfseries Catatan pengembalian} dari Administrator;
  \item {\bfseries Tanda tangan digital} Petugas BMN; dan
  \item {\bfseries Perangko digital}.
\end{itemize}

\subsection{Perangko Digital dan Tanda Tangan Digital}
\index{perangko digital}
\index{digital stamp}
\index{tanda tangan digital}
\index{digital signature}
\index{stempel}
Setelah borrower mengunggah Surat Pernyataan yang telah
ditandatangani dan distempel basah secara fisik, Administrator
dapat menambahkan {\bfseries perangko digital} dan
{\bfseries tanda tangan digital} ke dokumen. Proses ini
dilakukan secara terprogram menggunakan pustaka {\it pdf-lib}
yang memanipulasi dokumen PDF langsung di {\it server-side}.

\section{QR Code Verification}
\index{qr code verification}
\index{qr verification}
\index{qr scan}

Setiap transaksi yang disetujui dibangkitankan QR Code berisi
tautan verifikasi ke halaman detail transaksi. QR Code:
\begin{itemize}
  \item Dihasilkan dalam warna biru tua (\#1e3a5f) pada latar putih;
  \item Diunggah ke Vercel Blob dan disimpan URL-nya
        di basis data; dan
  \item Dapat dipindai menggunakan kamera atau dimasukkan secara
        manual oleh Administrator.
\end{itemize}

Hasil pemindaian langsung mengarahkan Administrator ke halaman
detail transaksi untuk mempercepat proses verifikasi dan
pengembalian.

\section{Impor dan Ekspor Data}
\index{impor data}
\index{ekspor data}
\index{excel}
\index{xlsx}

\subsection{Impor Barang}
\index{impor barang}
\index{import barang}
Importer barang dari file Excel mendukung strategi
{\it mirror sync}:
\begin{itemize}
  \item Barang baru {\bfseries ditambahkan};
  \item Barang yang ada {\bfseries diperbarui} informasinya;
  \item Barang yang tidak ada di file Excel {\bfseries dihapus}
        (dengan validasi tidak sedang dalam peminjaman aktif).
\end{itemize}

Importer juga mendukung field \texttt{tipe} yang baru ditambahkan
ke skema.

\subsection{Impor Pegawai}
\index{impor pegawai}
\index{import pegawai}
Importer data pegawai ({\it employee master}) dari Excel
untuk {\it seeding} atau pembaruan data akun Peminjam.

\subsection{Impor Peminjam + Pinjaman}
\index{impor peminjaman}
\index{import loan}
Importer untuk memigrasikan data dari sistem lama
({\it legacy system}) secara {\it batch}.

\subsection{Ekspor}
\index{ekspor}
\index{export}
Ekspor data ke format Excel tersedia untuk:
\begin{itemize}
  \item {\bfseries Barang}: seluruh data barang di lingkup akses;
  \item {\bfseries Peminjaman}: seluruh transaksi peminjaman
        dengan {\it filter} status; dan
  \item {\bfseries Pengguna}: seluruh data akun pengguna.
\end{itemize}

\section{Sistem Notifikasi}
\index{sistem notifikasi}
\index{notifications}
\index{in-app notification}

Notifikasi dalam aplikasi ({\it in-app notifications}) mencakup:
\begin{itemize}
  \item {\bfseries Pengajuan baru}: notifikasi ke Administrator
        saat borrower submits pengajuan;
  \item {\bfseries Status perubahan}: notifikasi ke borrower
        saat pengajuan disetujui, ditolak, atau diminta
        pengembalian;
  \item {\bfseries {\it Reminder} pengembalian}: notifikasi
        pengingat beberapa hari sebelum tanggal kembali
        rencana; dan
  \item {\bfseries Notifikasi pensiun}: Cron job harian
        yang mengirimkan notifikasi kepada Administrator
        dan borrower yang akan pensiun dalam 90 hari ke depan.
\end{itemize}

Notifikasi dapat ditandai sebagai dibaca, dihapus, dan
dapat diklik untuk navigasi langsung ke halaman terkait.

\section{Log Audit}
\index{audit log}
\index{activity log}
\index{audit trail}
\index{history}

Seluruh aktivitas penting dalam sistem dicatat ke tabel
{\it AuditLog}, meliputi:
\begin{itemize}
  \item Login dan logout;
  \item Pembuatan, pembaruan, dan penghapusan data;
  \item Perubahan status peminjaman;
  \item Impor dan ekspor data;
  \item Perubahan kata sandi; dan
  \item Pengalihan peran.
\end{itemize}

Setiap {\it log entry} menyimpan:
\begin{itemize}
  \item Timestamp kejadian;
  \item {\it User ID} yang melakukan;
  \item {\it Action} yang dilakukan;
  \item {\it Entity} dan {\it Entity ID} yang terlibat;
  \item {\it Snapshot} data sebelum perubahan (\texttt{dataLama})
        dan setelah perubahan (\texttt{dataBaru}) dalam format JSON;
  \item {\it IP address} sumber; dan
  \item {\it User agent} browser.
\end{itemize}

Pendekatan {\it denormalized snapshot} (tanpa {\it JOIN})
dipilih untuk menjamin keakuratan data historis meskipun
data utama kemudian diubah atau dihapus.

% ─── Chapter 13: Penjadwalan Otomatis ─────────────────────────────────────
\chapter{Penjadwalan Otomatis (Cron Jobs)}

\index{cron job}
\index{scheduled task}
\index{node-cron}
\index{automated task}

Sistem menjalankan beberapa tugas terjadwal secara otomatis
menggunakan pustaka {\it node-cron}:

\begin{longtable}[!htbp]{p{4cm} p{3cm} p{5.5cm}}
  \caption{Tugas Terjadwal}
  \label{tab:cron-jobs} \\
  \toprule
  \textbf{Tugas} & \textbf{Jadwal} & \textbf{Deskripsi} \\
  \midrule
  \endfirsthead
  \caption[]{Tugas Terjadwal (sambungan)} \\
  \toprule
  \textbf{Tugas} & \textbf{Jadwal} & \textbf{Deskripsi} \\
  \midrule
  \endhead
  \bottomrule
  \endfoot

  Notifikasi Pensiun       & Harian (00:00) & Memeriksa seluruh
                            pengguna dan mengirim notifikasi kepada
                            Administrator dan borrower yang akan
                            pensiun dalam 90 hari ke depan. Tanggal
                            pensiun dihitung secara otomatis dari NIP.
                            \\[6pt]
  Perhitungan Stok        & Saat diperlukan & Pengurangan dan
                            pemulihan stok barang dilakukan secara
                            atomik setiap kali status peminjaman
                            berubah (persetujuan, handover, atau
                            pengembalian). \\[6pt]
  Deteksi Terlambat        & Per request    & Setiap kali halaman
                            dashboard dimuat, sistem memeriksa
                            apakah ada peminjaman aktif yang telah
                            melewati batas tanggal kembali.
                            {\it Throttled} maksimal sekali per menit
                            per pemuatan. \\[6pt]
  Hapus Token Kadaluarsa   & Periodik        & {\it Cleanup} otomatis
                            untuk token yang telah kadaluarsa
                            pada tabel \texttt{BlacklistedToken}. \\
\end{longtable}

% ─── Chapter 14: Penutup ─────────────────────────────────────────────────────
\chapter{Penutup}

\section{Ringkasan}
\index{ringkasan}
\index{summary}

SIPP-BMN ({\it Sistem Informasi Peminjaman \& Pengembalian Barang
Milik Negara}) adalah aplikasi berbasis web yang komprehensif
untuk mengelola seluruh siklus hidup peminjaman dan pengembalian
Barang Milik Negara. Sistem ini menyediakan tiga peran utama:

\begin{enumerate}
  \item {\roleSuperAdmin} --- Akses penuh lintas Satker,
        pengelolaan admin, seluruh barang, seluruh transaksi,
        seluruh pengguna, seluruh Satker, dan log audit.

  \item {\roleAdmin} --- Pengelolaan barang di Satker tugas,
        persetujuan dan penolakan peminjaman, penandoveran
        dan pengembalian barang, pengelolaan peminjam di Satker,
        serta pemindaian QR Code.

  \item {\rolePeminjam} --- Akses katalog barang, pengajuan
        peminjaman (termasuk mode konsep), upload surat
        pernyataan, permintaan pengembalian, dan pelacakan
        riwayat pribadi.
\end{enumerate}

Fitur utama meliputi generasi dokumen otomatis (Surat Pernyataan
Peminjaman dan Surat Pengembalian) dengan perangko digital dan
tanda tangan digital, verifikasi QR Code, impor/ekspor data Excel,
pencegahan pemesanan ganda ({\it double-booking prevention}),
sistem notifikasi dalam aplikasi, log audit lengkap, penjadwalan
otomatis, serta mekanisme keamanan berlapis (JWT, Refresh Token,
CSRF Protection, Token Versioning, Rate Limiting, dan Security
Headers).

\section{Pengembang}
\index{pengembang}
\index{developer}
\index{credits}

Dokumentasi ini disusun berdasarkan kode sumber dan skema
basis data aplikasi SIPP-BMN. Untuk informasi lebih lanjut,
silakan merujuk pada dokumentasi teknis yang tersedia di
repositori proyek.

\begin{flushright}
  \textit{Tim Pengembang SIPP-BMN} \\[4pt]
  \today
\end{flushright}

% ─── Bibliography ───────────────────────────────────────────────────────────
\clearpage
\addcontentsline{toc}{chapter}{Daftar Pustaka}
\markboth{Daftar Pustaka}{Daftar Pustaka}

\begin{thebibliography}{9}
  \bibitem{nextjs}
    Vercel. \textit{Next.js Documentation}.
    \url{https://nextjs.org/docs}, 2024.

  \bibitem{prisma}
    Prisma. \textit{Prisma ORM Documentation}.
    \url{https://www.prisma.io/docs}, 2024.

  \bibitem{express}
    Express.js Team. \textit{Express.js Guide}.
    \url{https://expressjs.com}, 2024.

  \bibitem{jwt}
    IETF. \textit{JSON Web Token (JWT)}.
    RFC 7519. \url{https://tools.ietf.org/html/rfc7519}, 2015.

  \bibitem{rbac}
    NIST. \textit{Role-Based Access Control (RBAC)}.
    \url{https://csrc.nist.gov/projects/role-based-access-control}, 2024.

  \bibitem{pdf-lib}
    PDF-lib. \textit{PDF-lib: Create and modify PDF documents in the browser and Node.js}.
    \url{https://pdf-lib.js.org}, 2024.

  \bibitem{zod}
    Zod Team. \textit{Zod: TypeScript-first schema validation}.
    \url{https://zod.dev}, 2024.

  \bibitem{tailwind}
    Tailwind Labs. \textit{Tailwind CSS Documentation}.
    \url{https://tailwindcss.com/docs}, 2024.
\end{thebibliography}

% ─── Appendix A: Glossary ────────────────────────────────────────────────────
\clearpage
\addcontentsline{toc}{chapter}{Lampiran A --- Daftar Singkatan dan Istilah}
\markboth{Lampiran A: Daftar Singkatan dan Istilah}{Lampiran A: Daftar Singkatan dan Istilah}

\chapter*{Lampiran A: Daftar Singkatan dan Istilah}
\index{singkatan}
\index{glossary}

\begin{longtable}[!htbp]{p{3.5cm} p{9cm}}
  \caption{Daftar Singkatan dan Istilah}
  \label{tab:glossary} \\

  \toprule
  \textbf{Singkatan/Istilah} & \textbf{Definisi} \\
  \midrule
  \endfirsthead

  \caption[]{Daftar Singkatan dan Istilah (sambungan)} \\
  \toprule
  \textbf{Singkatan/Istilah} & \textbf{Definisi} \\
  \midrule
  \endhead

  \bottomrule
  \endfoot

  API        & Application Programming Interface \\
  BMN        & Barang Milik Negara (State-Owned Goods) \\
  CSRF       & Cross-Site Request Forgery \\
  CRUD       & Create, Read, Update, Delete \\
  ERD        & Entity Relationship Diagram \\
  FK         & Foreign Key (Kunci Tamu) \\
  HMN        & Barang Milik Negara (lihat BMN) \\
  HTTP       & Hypertext Transfer Protocol \\
  JWT        & JSON Web Token \\
  NIP        & Nomor Induk Pegawai (Employee Identification Number) \\
  NUP        & Nomor Urut (Item Serial Number) \\
  ORM        & Object-Relational Mapping \\
  RBAC       & Role-Based Access Control \\
  REST       & Representational State Transfer \\
  Satker     & Satuan Kerja (Work Unit / Cost Center) \\
  SIPP-BMN   & Sistem Informasi Peminjaman \& Pengembalian BMN \\
  SQL        & Structured Query Language \\
  UAT        & User Acceptance Testing \\
  URL        & Uniform Resource Locator \\
  XSS        & Cross-Site Scripting \\
\end{longtable}

% ─── Appendix B: Diagram ────────────────────────────────────────────────────
\clearpage
\addcontentsline{toc}{chapter}{Lampiran B --- Diagram Tambahan}
\markboth{Lampiran B: Diagram Tambahan}{Lampiran B: Diagram Tambahan}

\chapter*{Lampiran B: Diagram Tambahan}

\section*{Diagram Alur Pengajuan Peminjaman}
\index{diagram pengajuan}

\begin{figure}[htbp]
  \centering
  \IfFileExists{./assets/flow-pengajuan.png}
    {\includegraphics[width=0.8\textwidth]{./assets/flow-pengajuan.png}}
    {%
      \begin{tabular}{|p{3.5cm}|p{3.5cm}|p{3cm}|}
        \hline
        \textbf{Aktor} & \textbf{Tahap} & \textbf{Status} \\
        \hline
        Borrower & Pilih barang & --- \\
        Borrower & Tambah ke Keranjang & --- \\
        Borrower & Isi Formulir Ajukan & --- \\
        Borrower & Pilih Mode & Draft / Lengkap \\
        Borrower & Submit (mode Lengkap) & MENUNGGU \\
        Admin & Tinjauan & MENUNGGU \\
        Admin & Setuju & DISETUJUI \\
        Admin & Tolak & DITOLAK \\
        \hline
      \end{tabular}
    }
  \caption{Alur Pengajuan dan Persetujuan Peminjaman}
  \label{fig:flow-pengajuan}
\end{figure}

\section*{Diagram Alur Pengembalian}
\index{diagram pengembalian}

\begin{figure}[htbp]
  \centering
  \IfFileExists{./assets/flow-pengembalian.png}
    {\includegraphics[width=0.8\textwidth]{./assets/flow-pengembalian.png}}
    {%
      \begin{tabular}{|p{3.5cm}|p{3.5cm}|p{3cm}|}
        \hline
        \textbf{Aktor} & \textbf{Aksi} & \textbf{Hasil} \\
        \hline
        Borrower & --- & DIPINJAM / TERLAMBAT \\
        Borrower & Upload Surat Pengembalian & --- \\
        Admin & Terima notifikasi & --- \\
        Admin & Verifikasi barang & --- \\
        Admin & Konfirmasi Pengembalian & --- \\
        Borrower & --- & DIKEMBALIKAN (stok dipulihkan) \\
        Borrower & Berikan review & Selesai \\
        \hline
      \end{tabular}
    }
  \caption{Alur Pengembalian Barang}
  \label{fig:flow-pengembalian}
\end{figure}

% ─── Appendix C: Kode Status HTTP ───────────────────────────────────────────
\clearpage
\addcontentsline{toc}{chapter}{Lampiran C --- Kode Status HTTP}
\markboth{Lampiran C: Kode Status HTTP}{Lampiran C: Kode Status HTTP}

\chapter*{Lampiran C: Kode Status HTTP yang Digunakan}
\index{kode status HTTP}
\index{HTTP status codes}

\begin{longtable}[!htbp]{p{3cm} p{4cm} p{5.5cm}}
  \caption{Kode Status HTTP API}
  \label{tab:http-status} \\
  \toprule
  \textbf{Kode} & \textbf{Status} & \textbf{Penggunaan} \\
  \midrule
  \endfirsthead
  \caption[]{Kode Status HTTP API (sambungan)} \\
  \toprule
  \textbf{Kode} & \textbf{Status} & \textbf{Penggunaan} \\
  \midrule
  \endhead
  \bottomrule
  \endfoot

  200 & OK           & Request berhasil, data dikembalikan \\
  201 & Created      & Data berhasil dibuat \\
  204 & No Content   & Request berhasil, tidak ada body \\
  400 & Bad Request  & Input tidak valid atau request malformed \\
  401 & Unauthorized & Token tidak valid atau tidak ada \\
  403 & Forbidden    & Pengguna tidak memiliki izin \\
  404 & Not Found    & Data yang diminta tidak ditemukan \\
  409 & Conflict     & Konflik data (misalnya: double-booking) \\
  422 & Unprocessable Entity & Validasi gagal \\
  429 & Too Many Requests & Rate limit tercapai \\
  500 & Internal Server Error & Kesalahan server \\
\end{longtable}

% ─── Back Matter ────────────────────────────────────────────────────────────
\clearpage
\thispagestyle{empty}
\vspace*{2cm}
\begin{center}
  \HRule\\[0.3cm]
  {\LARGE\textbf{\textls{--- Akhir Dokumen ---}}}\\[0.3cm]
  \HRule\\[1cm]
  \begin{tabular}{c}
    \textbf{SIPP-BMN} \\
    \textit{Sistem Informasi Peminjaman \& Pengembalian} \\
    \textit{Barang Milik Negara} \\
    Versi 1.0.0
  \end{tabular}
\end{center}

\end{document}
